% Section 6.3, from lectures/L06/README.md: the objectives, the questions,
% the further reading, and what the next chapter does with all of it.


\section{Review}
\label{c6:sec:review}

\needspace{6\baselineskip}
\subsection*{What you should be able to do now}
After this chapter, you should be able to:
\begin{itemize}
  \item Find the quiescent point of a divider-biased stage, including the base current's effect.
  \item Say what divider stiffness is, choose it, and say what the usual rule of thumb costs.
  \item State the thermal drift of a base-emitter voltage, and the collector current drift it implies at a fixed base voltage.
  \item Explain what an emitter resistor does, in terms of feedback rather than in terms of a story about voltages moving.
  \item Compute the drift suppression, and notice it is the same number as the gain reduction.
  \item Apply the 220 mV rule and say where it comes from.
  \item Recognise an argument that is self-contradictory, from its own numbers.
\end{itemize}

\needspace{6\baselineskip}
\subsection*{Questions to test yourself}
You should be able to answer:
\begin{enumerate}
  \item A divider is designed for 1.71 V and measures 1.60 V with the transistor connected. Nothing is broken. What is happening, and which Chapter~\ref{c1:ch:circuits} result predicts it?
  \item Divider current ten times base current is the usual rule. What error in the collector current does that rule accept?
  \item At a fixed base voltage the collector current rises 8 per cent per degree. Where does that number come from, and which of $V_T$ or $I_S$ contributes most of it?
  \item An emitter resistor reduces drift and reduces gain. By what factor each, and why is that not a coincidence?
  \item The 220 mV rule puts about a fifth of a volt across the emitter resistor. What would 22 mV give, and what would 2.2 V give?
  \item Somebody argues that when the current rises 10 per cent, the base-emitter voltage falls from 0.65 V to 0.55 V and pulls it back. Find the contradiction.
  \item Your stage is biased at 1 mA and drifts 0.2 per cent per degree. What would it drift at 10 mA with the same emitter resistor?
\end{enumerate}

\needspace{6\baselineskip}
\subsection*{Further reading}
\begin{itemize}
  \item \emph{The Art of Electronics}, Horowitz and Hill, chapter 2, sections on biasing, for a second treatment of the same material.
\end{itemize}

\needspace{6\baselineskip}
\subsection*{Looking ahead}
\begin{itemize}
  \item Chapter~\ref{c7:ch:smallsignal} is the centre of the book. The small-signal model, the intrinsic emitter resistance the whole of Part~2 is written in terms of, and the emitter factor, including the node it actually belongs to.
\end{itemize}
